\documentclass[12pt]{article}
\usepackage[utf8]{inputenc}
\usepackage[T1]{fontenc}
\usepackage{hyperref}
\title{Wiarumus Language Resources for K-12 Curriculum Development}
\author{Compiled by Claude Sonnet 5 (claude-sonnet-5) for Andr\'as Kornai}
\date{}
\begin{document}
\maketitle

\section{Introduction}
Wiarumus (also called Mandi; ISO 639-3: tua) is a Torricelli language of Mandi village, Turubu Rural LLG, East Sepik Province, Papua New Guinea. The ethnic population is about 460 (2000 census). Only villagers born before about 1940 are reported to speak it, so it is severely endangered.

\section{Primary Educational Resources}
None were found.

\section{Linguistic and Reference Materials}
\begin{itemize}
\item Overview: \url{https://en.wikipedia.org/wiki/Wiarumus_language}
\item ASJP wordlist: \url{https://asjp.clld.org/languages/WIARUMUS}
\item Endangered Languages Project bibliography: \url{https://endangeredlanguages.com/lang/3489/bibliography}
\end{itemize}
No grammar or dictionary was found.

\section{Religious and Community Resources}
None were found.

\section{Digital and Audio Resources}
PARADISEC reportedly holds Wiarumus material in the Don Laycock (DL2) and William Foley (WF3) collections. The search results did not include direct item URLs, so none are cited.

\section{Curriculum Development Context}
Very little is documented and there are few remaining speakers. Any programme would begin with archival work (PARADISEC) and fieldwork with the elders.

\end{document}